\documentclass[11pt,a4paper]{article}
\usepackage[a4paper,margin=2.15cm]{geometry}
\usepackage[T1]{fontenc}
\usepackage[utf8]{inputenc}
\usepackage{lmodern,microtype}
\usepackage{amsmath,amssymb,mathtools,bm}
\usepackage{booktabs,enumitem,xcolor,hyperref,fancyhdr,array,needspace}

\definecolor{accent}{RGB}{35,70,120}
\hypersetup{colorlinks=true,linkcolor=accent,urlcolor=accent,citecolor=accent,
  pdftitle={Support-Function Harmonics for AR Sliding-Window Geometry: Convex-Hull Empirical Analogues},
  pdfauthor={Working notes for the Master's thesis}}
\setlist[itemize]{leftmargin=1.5em,itemsep=2pt,topsep=3pt}
\setlist[enumerate]{leftmargin=1.7em,itemsep=3pt,topsep=3pt}
\setlength{\parindent}{0pt}
\setlength{\parskip}{4.5pt}
\setlength{\headheight}{14pt}
\pagestyle{fancy}
\fancyhf{}
\lhead{Support-function harmonics for AR geometry}
\rhead{Convex-hull analogues}
\cfoot{\thepage}
\newcommand{\Sph}{\mathbb S}
\newcommand{\R}{\mathbb R}
\newcommand{\norm}[1]{\left\lVert #1\right\rVert}
\newcommand{\ip}[2]{\left\langle #1,#2\right\rangle}
\newcommand{\tr}{\operatorname{tr}}
\newcommand{\conv}{\operatorname{conv}}
\newcommand{\off}{\operatorname{off}}
\newcommand{\Frob}{\mathrm F}
\newcommand{\he}{\widehat h^{\mathrm e}}
\newcommand{\hh}{\widehat h}
\newcommand{\hQ}{\widehat Q}
\newcommand{\hC}{\widehat C}
\newcommand{\hM}{\widehat M}
\newenvironment{keybox}{%
  \begin{center}\begin{minipage}{0.94\linewidth}\color{accent}\hrule\vspace{5pt}\color{black}%
}{%
  \vspace{5pt}\color{accent}\hrule\color{black}\end{minipage}\end{center}%
}

\title{\textbf{Support-Function Harmonics for\\AR Sliding-Window Geometry}\\[5pt]
\large Bounded-support theory and convex-hull empirical analogues}
\author{Working notes for the Master's thesis}
\date{9 September 2026 \quad | \quad Revised version 6}

\begin{document}
\maketitle
\vspace{-1em}

\begin{keybox}
\textbf{Central construction.} For bounded innovations, the population sliding-window
support is a zonoid. Its empirical replacement throughout these notes is the ordinary
convex hull of the observed sliding-window point cloud:
\[
 \mathcal Z^k\quad\longmapsto\quad
 \widehat K_{k,T}=\conv\{Y_t^{(k)}:t=k,\ldots,T\},\qquad
 \hh_{k,T}(u)=\max_{k\le t\le T}u^\top Y_t^{(k)}.
\]
Projecting this support function onto spherical harmonics gives empirical size and
shape descriptors. Degree zero gives half the mean width, degree two gives quadratic
anisotropy, and degree four captures finer directional shape. All empirical quantities
below are functionals of this same convex hull.
\end{keybox}

\section{AR setup and the bounded population support}\label{sec:setup}
Consider the stationary causal AR($p$) process
\[
 X_t=\phi_1X_{t-1}+\cdots+\phi_pX_{t-p}+Z_t,
 \qquad Z_t\overset{\mathrm{i.i.d.}}\sim U[-1,1].
\]
Assume $\Phi(z)=1-\sum_{r=1}^p\phi_r z^r$ has no zeros for $|z|\le1$.
The impulse-response coefficients satisfy
\[
 \psi_0=1,\qquad \psi_j=\sum_{r=1}^p\phi_r\psi_{j-r}\ (j\ge1),
 \qquad \psi_j=0\ (j<0),\qquad \sum_{j\ge0}|\psi_j|<\infty.
\]
For $1\le k\le T$, put
\[
 Y_t^{(k)}=(X_t,X_{t-1},\ldots,X_{t-k+1})^\top,\qquad
 g_j^{(k)}=(\psi_j,\psi_{j-1},\ldots,\psi_{j-k+1})^\top.
\]
Then $Y_t^{(k)}=\sum_{j\ge0}g_j^{(k)}Z_{t-j}$. Independence and the full
innovation support $[-1,1]$ yield the compact, origin-symmetric support body
\[
 \boxed{\quad
 \mathcal Z^k=\sum_{j=0}^{\infty}[-g_j^{(k)},g_j^{(k)}],\qquad
 h_k(u)=\sum_{j=0}^{\infty}|u^\top g_j^{(k)}|.
 \quad}
\]
The infinite Minkowski sum is understood as its Hausdorff limit; convergence follows
from $\sum_j\norm{g_j^{(k)}}<\infty$. The empirical hull lies in $\mathcal Z^k$
almost surely. The uniform model fixes the innovation half-width at one; replacing
$[-1,1]$ by $[-b,b]$, $b>0$, multiplies population support functions and their
harmonic components by $b$.

\newpage
\section{Harmonic components and finite-sample symmetry}\label{sec:harmonics}
For $k\ge2$, let $\sigma_{k-1}$ be normalized surface measure on $\Sph^{k-1}$,
with total mass one. Write $\mathcal H_n^{(k)}$ for the spherical harmonics of degree
$n$, and
\[
 N(k,n)=\binom{k+n-1}{n}-\binom{k+n-3}{n-2},
 \qquad N(k,0)=1,\quad N(k,1)=k.
\]
The second binomial term is zero for $n<2$. For $k=2$, $N(2,n)=2$ for $n\ge1$;
the expansion is the usual Fourier series in the angle of $u$.
For an orthonormal basis $\{Y_{n,m}\}_{m=1}^{N(k,n)}$, define
\[
 a_{n,m}(K)=\int_{\Sph^{k-1}}h_K(v)Y_{n,m}(v)\,d\sigma(v),\qquad
 Q_n[K](u)=\sum_{m=1}^{N(k,n)}a_{n,m}(K)Y_{n,m}(u).
\]
Thus $h_K=\sum_{n\ge0}Q_n[K]$ in $L^2(\sigma)$ and
$\norm{Q_n[K]}_{L^2}^2=\sum_m|a_{n,m}(K)|^2$. Each $Q_n[K]$ is a function;
its norm is a basis-independent scalar descriptor. We use
\[
 Q_n^{(k)}:=Q_n[\mathcal Z^k],\qquad
 \hQ_n^{(k,T)}:=Q_n[\widehat K_{k,T}].
\]

\subsection*{One projection formula for both bodies}
Let $P_n^{(k)}(1)=1$ be the normalized zonal polynomial: a normalized Gegenbauer
polynomial for $k\ge3$, and $P_n^{(2)}(\cos\theta)=\cos(n\theta)$ for $k=2$.
The addition theorem gives
\[
 \boxed{\quad
 Q_n[K](u)=N(k,n)\int_{\Sph^{k-1}}
 h_K(v)P_n^{(k)}(u^\top v)\,d\sigma(v).
 \quad}
\]
In particular,
\[
 P_2^{(k)}(t)=\frac{kt^2-1}{k-1},\qquad
 P_4^{(k)}(t)=\frac{(k+2)(k+4)t^4-6(k+2)t^2+3}{(k-1)(k+1)}.
\]

\subsection*{The raw hull and its even support component}
The population body is origin-symmetric, so $Q_{2r+1}^{(k)}=0$. A finite-sample
hull usually has nonzero odd components. Its even part is
\[
 \boxed{\quad
 \he_{k,T}(u)=\frac{\hh_{k,T}(u)+\hh_{k,T}(-u)}2
 =\frac{\max_tu^\top Y_t^{(k)}-\min_tu^\top Y_t^{(k)}}2.
 \quad}
\]
This is the support function of $\tfrac12(\widehat K_{k,T}-\widehat K_{k,T})$.
Using it in an even-degree projection leaves the raw-hull result unchanged:
\[
 \hQ_{2r}^{(k,T)}=\Pi_{2r}\he_{k,T},\qquad
 \he_{k,T}=\hQ_0^{(k,T)}+\hQ_2^{(k,T)}+\hQ_4^{(k,T)}+\cdots
 \quad\text{in }L^2.
\]
This is a way to compute the original hull's even harmonics and widths. It is
different from replacing the sample by $\conv\{\pm Y_t^{(k)}\}$, whose support
function is $\max_t|u^\top Y_t^{(k)}|$ and whose widths generally differ.

\newpage
\section{Bounded-support formulas for the theoretical target}\label{sec:population}
These formulas describe $\mathcal Z^k$ and provide a benchmark for the empirical
hull. Write $g_j^{(k)}=r_jv_j$, where $r_j=\norm{g_j^{(k)}}$ and $v_j=g_j^{(k)}/r_j$;
omit zero generators. The even generating measure is
\[
 \rho_k=\frac12\sum_{j\ge0}r_j(\delta_{v_j}+\delta_{-v_j}),\qquad
 h_k(u)=\int_{\Sph^{k-1}}|u^\top v|\,\rho_k(dv).
\]
The spherical cosine transform acts diagonally on each harmonic degree. Put
\[
 \mu_k=\frac{\Gamma(k/2)}{\sqrt\pi\,\Gamma((k+1)/2)},\qquad
 \lambda_{n,k}=\int_{\Sph^{k-1}}|v_1|P_n^{(k)}(v_1)\,d\sigma(v).
\]
Then
\[
 \lambda_{0,k}=\mu_k,\qquad
 \lambda_{2,k}=\frac{\mu_k}{k+1},\qquad
 \lambda_{4,k}=-\frac{\mu_k}{(k+1)(k+3)},
\]
and $Q_n^{(k)}(u)=N(k,n)\lambda_{n,k}\sum_jr_jP_n^{(k)}(u^\top v_j)$.

\subsection*{Degree zero: mean width}
With $S_k=\sum_jr_j$,
\[
 \boxed{\quad Q_0^{(k)}=\mu_k S_k,\qquad
 \overline w(\mathcal Z^k)=2Q_0^{(k)}=2\mu_k\sum_j\norm{g_j^{(k)}}.\quad}
\]
Equivalently, $\mu_k=2\kappa_{k-1}/(k\kappa_k)$, where $\kappa_k$ is the volume
of the unit ball in $\R^k$. Since $V_1(\mathcal Z^k)=2\sum_jr_j$,
the mean width equals $\mu_k V_1(\mathcal Z^k)$.

\subsection*{Degree two: quadratic anisotropy}
Define
\[
 M_k=\sum_{j\ge0}\frac{g_j^{(k)}g_j^{(k)\top}}{\norm{g_j^{(k)}}},\qquad
 A_k=M_k-\frac{S_k}{k}I_k,\qquad
 c_k=\frac{\mu_k k(k+2)}{2(k+1)}.
\]
Then $\tr M_k=S_k$, $\tr A_k=0$, and
\[
 \boxed{\quad Q_2^{(k)}(u)=c_k u^\top A_ku,\qquad
 \norm{Q_2^{(k)}}_{L^2}^2=
 \mu_k^2\frac{k(k+2)}{2(k+1)^2}\norm{A_k}_{\Frob}^2.\quad}
\]
The matrix $C_k:=c_kA_k$ is the direct quadratic harmonic coefficient matrix.

\subsection*{Degree four: finer directional shape}
With $N(k,4)=k(k-1)(k+1)(k+6)/24$,
\[
 Q_4^{(k)}(u)=N(k,4)\lambda_{4,k}\sum_jr_jP_4^{(k)}(u^\top v_j).
\]
For any degree $n$, the corresponding generator identity is
\[
 \norm{Q_n^{(k)}}_{L^2}^2
 =N(k,n)\lambda_{n,k}^2\sum_{i,j\ge0}r_ir_jP_n^{(k)}(v_i^\top v_j).
\]
These generator formulas are specific to the population support representation.
The empirical analogues below instead insert the hull support function into the
general harmonic projection formula.

\newpage
\section{Direct empirical analogues of \texorpdfstring{$Q_0$ and $Q_2$}{Q0 and Q2}}\label{sec:empirical}
For readability, suppress $(k,T)$ within this section and write
$\hh=h_{\widehat K_{k,T}}$, $\he=(\hh(\cdot)+\hh(-\cdot))/2$.
All displayed integrals are exact functionals of the observed hull.

\subsection*{Empirical mean width}
The degree-zero component is constant and satisfies
\[
 \boxed{\quad
 \hQ_0=\int_{\Sph^{k-1}}\hh(u)\,d\sigma(u)
 =\int_{\Sph^{k-1}}\he(u)\,d\sigma(u),\qquad
 \widehat{\overline w}_{k,T}=2\hQ_0.
 \quad}
\]
Consequently,
\[
 \widehat{\overline w}_{k,T}
 =\int_{\Sph^{k-1}}\left(\max_tu^\top Y_t^{(k)}-
 \min_tu^\top Y_t^{(k)}\right)\,d\sigma(u).
\]
This is the spherical average of the observed projection range. It depends only on
the convex hull, so adding points in its interior does not change the result.

\subsection*{The traceless quadratic coefficient matrix}
Define
\[
 \boxed{\quad
 \hC_{k,T}=\frac{k(k+2)}2\int_{\Sph^{k-1}}
 \he(u)\left(uu^\top-\frac1kI_k\right)\,d\sigma(u).
 \quad}
\]
The same matrix results if $\he$ is replaced by $\hh$.
Using the normalized spherical fourth moments gives
\[
 \boxed{\quad
 \hQ_2^{(k,T)}(u)=u^\top\hC_{k,T}u,\qquad
 \norm{\hQ_2^{(k,T)}}_{L^2}^2=
 \frac{2}{k(k+2)}\norm{\hC_{k,T}}_{\Frob}^2.
 \quad}
\]
The matrix is symmetric and traceless. Its eigenvectors describe the principal
directions of the degree-two component, while its eigenvalues quantify signed
deviations from the degree-zero baseline. In particular,
$\int\hQ_2^{(k,T)}\,d\sigma=0$.

\subsection*{Recovering an empirical counterpart of the support matrix $M_k$}
To connect the hull harmonics to the population matrix in Section~\ref{sec:population}, set
\[
 \widehat A_{k,T}^{\mathrm{sup}}=\frac{\hC_{k,T}}{c_k},\qquad
 \boxed{\quad
 \hM_{k,T}^{\mathrm{sup}}=
 \frac{\hC_{k,T}}{c_k}+\frac{\hQ_0^{(k,T)}}{k\mu_k}I_k.
 \quad}
\]
Equivalently, if $J_{k,T}=\int\he(u)uu^\top\,d\sigma(u)$, then
\[
 \boxed{\quad
 \hM_{k,T}^{\mathrm{sup}}=
 \frac1{\mu_k}\left((k+1)J_{k,T}-\hQ_0^{(k,T)}I_k\right),\qquad
 \tr\hM_{k,T}^{\mathrm{sup}}=\frac{\hQ_0^{(k,T)}}{\mu_k}.
 \quad}
\]
This is a linear reconstruction from the hull's degree-zero and degree-two
harmonics. A general empirical hull need not be a zonoid, so the reconstruction
does not supply a generator representation of the entire hull.
It targets $M_k$ when the empirical hull converges to $\mathcal Z^k$.

\newpage
\section{Empirical \texorpdfstring{$Q_4$}{Q4} and numerical spherical integration}\label{sec:numerics}
The direct degree-four analogue is
\[
 \boxed{\quad
 \hQ_4^{(k,T)}(u)=N(k,4)\int_{\Sph^{k-1}}
 \he_{k,T}(v)P_4^{(k)}(u^\top v)\,d\sigma(v).
 \quad}
\]
The addition theorem also gives the exact, basis-free energy identity
\[
 \boxed{\quad
 \norm{\hQ_4^{(k,T)}}_{L^2}^2
 =N(k,4)\iint
 \he_{k,T}(v)\he_{k,T}(w)P_4^{(k)}(v^\top w)\,d\sigma(v)d\sigma(w).
 \quad}
\]
Both integrations are over the unit sphere. These expressions use directional
maxima and minima of the point cloud; population generator multipliers are not
inserted into this general projection formula.

\subsection*{A shared set of integration directions}
Choose $L$ unit directions $U_\ell$ with nonnegative weights $\omega_\ell$ summing
to one. Independent uniform directions with $\omega_\ell=1/L$ provide a Monte
Carlo rule. In dimension two, equally spaced angles give a simple deterministic
rule. For higher dimensions, spherical designs or other refined cubature rules
can be used. Evaluate
\[
 H_\ell=\frac12\left(\max_tU_\ell^\top Y_t^{(k)}-
 \min_tU_\ell^\top Y_t^{(k)}\right),\qquad \ell=1,\ldots,L.
\]
Then the numerical approximations are
\begin{align*}
 \hQ_0^{(k,T,L)}&=\sum_{\ell=1}^L\omega_\ell H_\ell,\\
 \hC_{k,T}^{(L)}&=\frac{k(k+2)}2\sum_{\ell=1}^L\omega_\ell H_\ell
 \left(U_\ell U_\ell^\top-\frac1kI_k\right),\\
 \hQ_4^{(k,T,L)}(u)&=N(k,4)\sum_{\ell=1}^L\omega_\ell H_\ell
 P_4^{(k)}(u^\top U_\ell).
\end{align*}
Once these harmonic approximations have been formed, their spherical norms satisfy
\begin{align*}
 \norm{\hQ_2^{(k,T,L)}}_{L^2}^2
 &=\frac{2}{k(k+2)}\norm{\hC_{k,T}^{(L)}}_{\Frob}^2,\\
 \norm{\hQ_4^{(k,T,L)}}_{L^2}^2
 &=N(k,4)\sum_{\ell,m=1}^L\omega_\ell\omega_mH_\ell H_m
 P_4^{(k)}(U_\ell^\top U_m).
\end{align*}
The last expression is nonnegative in exact arithmetic and is the exact norm of
the numerically constructed degree-four function. At finite $L$ it is an
approximation to the exact hull energy; under Monte Carlo integration it is
generally not an unbiased estimator of that energy.

\textbf{Accuracy and cost.} Support evaluation requires only projection maxima and
minima, so constructing hull facets is unnecessary. Batched projections cost
$O(L(T-k+1)k)$; the direct degree-four pair sum costs $O(L^2k)$.
Because the hull support function is piecewise linear on directions, a low-degree
spherical design does not make these integrals exact. Increase $L$ and compare
independent direction sets, or refine the cubature, until the descriptors stabilize.

\newpage
\section{Minimum width, location, and shape summaries}\label{sec:width}
\subsection*{Exact empirical width and its degree-two approximation}
For the raw hull, width in direction $u$ and minimum width are
\[
 \widehat w_{k,T}(u)=\hh_{k,T}(u)+\hh_{k,T}(-u)=2\he_{k,T}(u),\qquad
 \widehat w_{\min,k,T}=2\min_{\norm u=1}\he_{k,T}(u).
\]
Keeping degrees zero and two gives
\[
 \boxed{\quad
 \widehat w_{\min,k,T}^{(2)}
 =2\left(\hQ_0^{(k,T)}+\lambda_{\min}(\hC_{k,T})\right).
 \quad}
\]
Thus $2\hQ_0$ is the mean-width baseline, and $2\lambda_{\min}(\hC)$ is the
leading correction toward the thinnest direction. The population counterpart is
$w_{\min}^{(2)}(k)=2(Q_0^{(k)}+c_k\lambda_{\min}(A_k))$.
For the continuous residual
\[
 \widehat R_{\ge4}(u)=\he_{k,T}(u)-\hQ_0^{(k,T)}-u^\top\hC_{k,T}u,
\]
the elementary inequality $|\min f-\min g|\le\norm{f-g}_\infty$ yields
\[
 \boxed{\quad
 |\widehat w_{\min,k,T}-\widehat w_{\min,k,T}^{(2)}|
 \le2\norm{\widehat R_{\ge4}}_\infty.
 \quad}
\]
The truncated function need not be a support function, and its minimum can be
negative. A small $L^2$ energy in degree four alone does not bound the full
uniform residual. Comparing the mean-width, degree-two minimum-width, and full
minimum-width curves tests how much of a transition is already explained by $Q_2$.

\textbf{Directional search.} The smallest width among finitely many sampled
directions is an upper bound on the exact empirical minimum width. It should be
reported as an approximation unless the minimization has been certified.
If the cloud has affine dimension below $k$, its exact minimum width is zero;
in particular, this occurs whenever $T-k+1\le k$.

\subsection*{Steiner point and odd harmonics}
The raw hull has the empirical Steiner point
\[
 \widehat s_{k,T}=k\int_{\Sph^{k-1}}u\,\hh_{k,T}(u)\,d\sigma(u),\qquad
 \hQ_1^{(k,T)}(u)=u^\top\widehat s_{k,T}.
\]
It belongs to the hull and is translation and rotation equivariant. It need not
vanish in a finite sample, whereas $s(\mathcal Z^k)=0$. Translating the hull by
$-\widehat s_{k,T}$ removes degree one; it does not generally remove higher odd
degrees. Even harmonics and widths are translation invariant. Their computation
from $\he_{k,T}$ automatically removes all odd contributions.

\subsection*{Scale-free empirical shape diagnostics}
When $\hQ_0^{(k,T)}>0$, define
\[
 \widehat R_2(k)=\frac{\norm{\hQ_2^{(k,T)}}_{L^2}}{\hQ_0^{(k,T)}},\qquad
 \widehat R_4(k)=\frac{\norm{\hQ_4^{(k,T)}}_{L^2}}{\hQ_0^{(k,T)}},\qquad
 \widehat D_2(k)=-\frac{\lambda_{\min}(\hC_{k,T})}{\hQ_0^{(k,T)}}.
\]
They summarize relative anisotropy. In particular,
$\widehat w_{\min,k,T}^{(2)}/\widehat{\overline w}_{k,T}=1-\widehat D_2(k)$.

\newpage
\section{Exact AR-filter structure in the bounded-support setting}\label{sec:ar}
Assume $p\ge1$ is the minimal AR order, so $\phi_p\ne0$. At dimension $p+1$ the
complete recursion first fits inside one sliding window. With
$a=(1,-\phi_1,\ldots,-\phi_p)^\top$ and $u_a=a/\norm a$,
\[
 a^\top Y_t^{(p+1)}=Z_t,\qquad
 \boxed{\quad h_{p+1}(u_a)=\frac1{\norm a},\qquad
 w_{\mathcal Z^{p+1}}(u_a)=\frac2{\norm a}.\quad}
\]
For the observed raw hull, the corresponding exact identities are
\[
 \hh_{p+1,T}(u_a)=\frac{\max_{t=p+1,\ldots,T}Z_t}{\norm a},\qquad
 \boxed{\quad
 \widehat w_{p+1,T}(u_a)=
 \frac{\max_{t=p+1,\ldots,T}Z_t-\min_{t=p+1,\ldots,T}Z_t}{\norm a}.
 \quad}
\]
This bounds the empirical minimum width from above. It does not imply that $u_a$
is the globally thinnest direction or an eigenvector of the degree-two matrix.

\subsection*{A growing family of filter directions}
For $k\ge p+1$, let $B_{p,k}\in\R^{(k-p)\times k}$ have shifted copies of $a^\top$
as its rows. Then
\[
 B_{p,k}Y_t^{(k)}=(Z_t,Z_{t-1},\ldots,Z_{t-k+p+1})^\top.
\]
Its row space $\mathcal S_{p,k}=\operatorname{range}(B_{p,k}^\top)$ has dimension
$k-p$. For $c\in\R^{k-p}\setminus\{0\}$ and
$u=B_{p,k}^\top c/\norm{B_{p,k}^\top c}$, independence and the cube support of
the innovation block give
\[
 h_k(u)=\frac{\norm{c}_1}{\norm{B_{p,k}^\top c}}.
\]
The appearance of $1,2,3,\ldots$ filter directions at dimensions
$p+1,p+2,p+3,\ldots$ is exact algebra. A separated eigenvalue block or a width
elbow requires additional evidence.

\subsection*{Congruence identity for the population support matrix}
The impulse-response recursion gives
\[
 B_{p,k}g_j^{(k)}=
 \begin{cases}
 e_{j+1},&0\le j\le k-p-1,\\
 0,&j\ge k-p,
 \end{cases}
\]
where $e_1,\ldots,e_{k-p}$ are standard basis vectors. Substitution into the
generator formula for $M_k$ yields
\[
 \boxed{\quad
 B_{p,k}M_kB_{p,k}^\top
 =\operatorname{diag}\left(
 \frac1{\norm{g_0^{(k)}}},\ldots,\frac1{\norm{g_{k-p-1}^{(k)}}}
 \right).
 \quad}
\]
This is a congruence identity; the filter matrix is not generally orthogonal.
It does not imply repeated eigenvalues. Replacing $M_k$ by
$\hM_{k,T}^{\mathrm{sup}}$ gives an empirical diagnostic whose off-diagonal
entries need not vanish at finite $T$.

\newpage
\section{Exploratory order selection from the empirical hull}\label{sec:order}
\subsection*{Primary candidate: approximate filter congruence}
For a candidate order $q$, let $\beta\in\R^q$ be its coefficient vector.
Form $B_{q,k}(\beta)$ from shifted copies of $(1,-\beta_1,\ldots,-\beta_q)$.
Define
\[
 \widehat D_{q,k,T}(\beta)=
 B_{q,k}(\beta)\hM_{k,T}^{\mathrm{sup}}B_{q,k}(\beta)^\top,
\]
and, on a specified admissible coefficient set $\Theta_q$,
\[
 \boxed{\quad
 \widehat R_{q,k,T}^{*}=
 \inf_{\substack{\beta\in\Theta_q\\\norm{\widehat D_{q,k,T}(\beta)}_{\Frob}>0}}
 \frac{\norm{\off(\widehat D_{q,k,T}(\beta))}_{\Frob}^2}
 {\norm{\widehat D_{q,k,T}(\beta)}_{\Frob}^2}.
 \quad}
\]
Here $\off$ sets the diagonal to zero. A vanishing denominator is excluded and
must be flagged if encountered numerically. For $q=0$, take $B_{0,k}=I_k$ with
no fitted coefficients. If orders $0,\ldots,q_{\max}$ are compared at one dimension,
use $k\ge q_{\max}+2$: a one-row filter gives a $1\times1$ matrix and therefore
an automatically zero off-diagonal score.

The population score $R_{q,k}^{*}$ is obtained by replacing
$\hM_{k,T}^{\mathrm{sup}}$ with $M_k$. At $q=p$ it is zero when $\Theta_p$
contains the true coefficients. Higher orders also attain zero when their
admissible sets contain the zero-padded true filter. Thus a candidate rule seeks
the \emph{smallest} order compatible with a near-zero score, or a substantial last
decrease in
\[
 \widehat\Delta_q(k)=\widehat R_{q-1,k,T}^{*}-\widehat R_{q,k,T}^{*}.
\]
Identifiability against $q<p$, threshold calibration, and reliability of numerical
optimization remain to be established. The exact population identity alone does
not establish an order estimator. Use the same sample windows and integration
directions across all candidate orders at a given $k$.

\subsection*{Secondary diagnostics: spectra and width curves}
Let $\widehat\gamma_1^{(k,T)}\le\cdots\le\widehat\gamma_k^{(k,T)}$ be the
eigenvalues of $\hC_{k,T}$. For $1\le r<k$ and $\norm{\hC_{k,T}}_{\Frob}>0$,
an exploratory gap statistic is
\[
 \widehat G_k(r)=
 \frac{\widehat\gamma_{k-r+1}^{(k,T)}-\widehat\gamma_{k-r}^{(k,T)}}
 {\norm{\hC_{k,T}}_{\Frob}}.
\]
Compare these gaps across $k$, and in simulations compare eigenvectors or
eigenspaces with the known filter directions. A low-eigenvalue subspace of
dimension $k-p$ is a hypothesis, not a consequence of the congruence identity.

The hull mean width $2\hQ_0^{(k,T)}$, the degree-two minimum-width approximation,
and the full minimum width provide complementary curves. An elbow or L-method
may be applied as an exploratory comparison, but no decreasing or flattening
pattern is assumed. In particular, a breakpoint at $k=p+1$ must be tested for
these hull quantities. Using a different width curve does not by itself resolve
the original breakpoint-selection problem.

\newpage
\section{Consistency and the two sources of approximation error}\label{sec:consistency}
Fix $k$. For a stationary ergodic window sequence whose marginal support is the
compact body $\mathcal Z^k$, the observed points are almost surely dense in that
support as $T\to\infty$. Indeed, every open neighborhood of a support point has
positive probability and is visited by ergodicity. A finite cover of the compact
support then gives
\[
 \delta_{k,T}:=d_H(\widehat K_{k,T},\mathcal Z^k)\longrightarrow0
 \quad\text{almost surely}.
\]
This applies to the stationary causal AR model in Section~\ref{sec:setup}; the
overlapping windows are not assumed independent. It is a fixed-dimension
statement and supplies no rate uniform in growing $k$.

\subsection*{Geometric error from the finite point cloud}
For compact convex bodies, $d_H(K,L)=\norm{h_K-h_L}_\infty$. Orthogonal projection
is contractive in $L^2(\sigma)$, hence, for every fixed harmonic degree $n$,
\[
 \boxed{\quad
 \norm{\hQ_n^{(k,T)}-Q_n^{(k)}}_{L^2}\le\delta_{k,T},\qquad
 \left|\norm{\hQ_n^{(k,T)}}_{L^2}-\norm{Q_n^{(k)}}_{L^2}\right|
 \le\delta_{k,T}.
 \quad}
\]
In particular,
\[
 |\hQ_0^{(k,T)}-Q_0^{(k)}|\le\delta_{k,T},\qquad
 \norm{\hC_{k,T}-C_k}_{\Frob}
 \le\sqrt{\frac{k(k+2)}2}\,\delta_{k,T}.
\]
The matrix reconstruction in Section~\ref{sec:empirical} therefore converges to
$M_k$. The exact widths satisfy
\[
 |\overline w(\widehat K_{k,T})-\overline w(\mathcal Z^k)|\le2\delta_{k,T},\qquad
 |\widehat w_{\min,k,T}-w_{\min}(\mathcal Z^k)|\le2\delta_{k,T}.
\]
Since both bodies are measured in the same coordinates, these are direct
comparisons. The degree-two width approximation is also consistent for its
population counterpart by continuity of the smallest eigenvalue. Its truncation
error relative to the full population minimum width is a separate issue.

\subsection*{Numerical error from finitely many directions}
At fixed $T$, independent uniform integration directions give convergence of the
numerical coefficients as $L\to\infty$. The practical error decomposes as
\[
 \norm{\hQ_n^{(k,T,L)}-Q_n^{(k)}}_{L^2}
 \le
 \underbrace{\norm{\hQ_n^{(k,T,L)}-\hQ_n^{(k,T)}}_{L^2}}
 _{\text{spherical integration error}}
 +\underbrace{\delta_{k,T}}_{\text{finite-cloud error}}.
\]
Increasing $L$ improves integration accuracy; increasing $T$ improves coverage
of the support. These should be varied separately in numerical experiments.

For the congruence criterion, compact admissible coefficient sets and denominators
uniformly bounded away from zero allow matrix convergence to pass to uniform
convergence of the objective and convergence of the infimum values. Consistency
of selected AR orders additionally requires identifiability and a justified
selection rule. Eigenspace convergence likewise requires a separating population
eigengap; matrix convergence alone is insufficient.

\newpage
\section{Focused workflow and interpretation}\label{sec:workflow}
\begin{enumerate}
\item \textbf{Fix the bounded AR model and coordinate scale.} Use a stationary
sample, or document simulation burn-in. Form the raw windows $Y_t^{(k)}$ for each
chosen $k$. Keep the coordinate order and any preprocessing explicit.
\item \textbf{Compute hull support values.} For each integration direction, record
the maximum and minimum projected window. Their half-range is $H_\ell$.
\item \textbf{Obtain the harmonic descriptors.} Compute $\hQ_0$, $\hC$,
$\norm{\hQ_2}_{L^2}$, $\norm{\hQ_4}_{L^2}$, and
$\hM^{\mathrm{sup}}$ with the formulas in Sections~\ref{sec:empirical}--\ref{sec:numerics}.
\item \textbf{Compare width and shape.} Plot mean width, the degree-two
minimum-width approximation, a clearly labeled numerical full-width estimate,
and the scale-free ratios. Check for affine degeneracy before interpreting
minimum width.
\item \textbf{Test the AR-specific criterion.} Evaluate the congruence scores for
candidate orders at dimensions with at least two filter rows. In simulations,
compare with the exact population generator calculation. Treat spectral gaps and
elbows as secondary diagnostics.
\item \textbf{Separate validation tasks.} Refine the directional integration at
fixed data, then vary sample size and repeat simulations across stable coefficient
regimes. Establish lower-order identifiability before claiming consistent order
recovery.
\end{enumerate}

For $k=1$, the hull is the interval $[\min_tX_t,\max_tX_t]$. Its mean and minimum
widths both equal the observed range, and $\hQ_0$ equals half that range. The
spherical degree-two and degree-four formulas in these notes apply only to $k\ge2$.

\begin{center}
\renewcommand{\arraystretch}{1.22}
\begin{tabular}{@{}p{3.35cm}p{12.9cm}@{}}
\toprule
\textbf{Object} & \textbf{Role in the empirical hull analysis}\\
\midrule
$\widehat K_{k,T}$ & Convex hull of the observed sliding-window point cloud.\\
$\hh_{k,T}$ & Directional maximum; determines the whole hull.\\
$\he_{k,T}$ & Half the directional range; preserves all even hull harmonics.\\
$\hQ_0^{(k,T)}$ & Half the empirical mean width.\\
$\hC_{k,T}$ & Traceless matrix encoding the empirical degree-two component.\\
$\widehat w_{\min,k,T}^{(2)}$ & Mean-width baseline plus the leading thin-direction correction.\\
$\hQ_4^{(k,T)}$ & Degree-four directional shape from the same hull.\\
$\hM_{k,T}^{\mathrm{sup}}$ & Reconstruction of the population support matrix from $\hQ_0+\hQ_2$.\\
$\widehat R_{q,k,T}^{*}$ & Proposed AR-order diagnostic based on approximate filter congruence.\\
\bottomrule
\end{tabular}
\end{center}

\subsection*{References and normalization}
H. Groemer, \emph{Geometric Applications of Fourier Series and Spherical Harmonics},
Cambridge University Press, 1996. Background on the addition theorem,
Funk--Hecke formula, and support-function harmonics.
\href{https://www.cambridge.org/core/books/geometric-applications-of-fourier-series-and-spherical-harmonics/13DF62702900BE0560420729C0808474}{Publisher reference}.

R. Schneider, \emph{Integral Geometric Tools for Stochastic Geometry}, lecture
notes. Background on support functions and the Steiner point.
\href{https://wwwalt.math.kit.edu/stoch/seite/konvgeo_cime04/media/notes_schneider.pdf}{Author's lecture notes}.

All spherical integrals and $L^2$ norms use $\sigma_{k-1}(\Sph^{k-1})=1$.
The AR-filter identities are derived above; order-selection proposals are stated
as research questions, rather than attributed to these background references.

\end{document}
